%%%%%%%%%%%%%%%%%%%%%%%%%%%%%%%%%%%%%%%%%%%%%%%%%%%%%%%%%%%
% Dead Code Detection — VerCors Implementation Report
%%%%%%%%%%%%%%%%%%%%%%%%%%%%%%%%%%%%%%%%%%%%%%%%%%%%%%%%%%%

\documentclass[9pt,a4paper,twocolumn]{article}
\usepackage[english]{babel}
\usepackage[T1]{fontenc}
\usepackage[utf8]{inputenc}
\usepackage{listings}
\usepackage{xcolor}
\usepackage{booktabs}
\usepackage{array}
\usepackage{geometry}
\usepackage{titlesec}
\usepackage{parskip}
\usepackage{microtype}
\usepackage{float}
\usepackage[hidelinks]{hyperref}

\geometry{a4paper, margin=1.8cm, columnsep=0.6cm}

\lstset{
  basicstyle=\ttfamily\small,
  frame=single,
  breaklines=true,
  columns=fullflexible,
  keepspaces=true,
  xleftmargin=4pt,
  xrightmargin=4pt
}

\titleformat{\section}{\normalfont\large\bfseries}{\thesection}{0.8em}{}
\titleformat{\subsection}{\normalfont\normalsize\bfseries}{\thesubsection}{0.8em}{}
\titleformat{\subsubsection}{\normalfont\normalsize\itshape}{\thesubsubsection}{0.8em}{}
\titlespacing{\section}{0pt}{8pt}{4pt}
\titlespacing{\subsection}{0pt}{6pt}{2pt}
\titlespacing{\subsubsection}{0pt}{4pt}{2pt}

\setlength{\parskip}{4pt}
\setlength{\parindent}{0pt}

\title{\textbf{Smoke Testing in VerCors}\\
\large Dead Code and Inconsistent Specification Detection}
\author{Alexandra Jakovleva}
\date{\today}

\begin{document}

\maketitle

%----------------------------------------------------------
\section{Introduction}

When a method's specification is contradictory, its branches can
never execute but the verifier proves everything inside them
valid anyway, without any warning. This is dead code
that looks verified but never actually runs. The problem is subtle:
the developer sees a green checkmark, but the verified code is
unreachable in any real execution.

This report documents \texttt{DetectDeadCode}, a rewrite
pass that automatically inserts \texttt{refute false} at strategic
points in the program, using VerCors itself without any external sources. If the verifier can prove \texttt{false}
at such a point, the code is unreachable and a \texttt{deadBranch}
warning is produced. A small addition is a
\texttt{CheckLoopInvariantSatisfiability} pass which handles the specific case
of self-contradictory loop invariants in an isolated proof context.

\textbf{Why \texttt{refute false}?} Frama-C inserts
\texttt{assert false} at smoke-test points. In VerCors, a direct
equivalent would require generating a separate helper procedure for
each check (similar to the existing
\texttt{CheckContractSatisfiability} pass) and passing all relevant
preconditions explicitly. This approach loses the mid-body symbolic
state: \texttt{assume} and \texttt{inhale} statements, function
calls, and value manipulations all contribute to the state but
cannot easily be extracted into a fresh procedure. Inserting
\texttt{refute false} directly inside the branch body avoids this
entirely: the verifier already holds the full symbolic state at
that point — preconditions, assumptions, inhaled facts, loop
invariants — with no extraction needed. The check passes silently
on live paths and fails only when the verifier can prove
\texttt{false} from the current state, which happens when
the code is unreachable.

\textbf{Pipeline placement.} The pass runs in the Transformation
stage after \texttt{PropagateContextEverywhere} and
\texttt{EncodeRangedFor}, and before
\texttt{RefuteToInvertedAssert} and \texttt{ParBlockEncoder}.
Language-specific lowering (\texttt{LangSpecificToCol}) has already
been completed in the earlier Resolution stage. The placement is
deliberate: by this point \texttt{for} loops are
already normalised to \texttt{Loop} nodes,
\texttt{context\_everywhere} clauses are merged into loop invariants,
and \texttt{ParBlock}/\texttt{ParAtomic} nodes still exist in the
AST and can be instrumented. The generated \texttt{Refute(false)}
nodes are not lowered by \texttt{RefuteToInvertedAssert}; instead
they pass through the pipeline unchanged and are translated
directly to Viper's native \texttt{refute} statement in
\texttt{ColToSilver} via the \texttt{RefutePlugin}.

The following sections walk through each feature of the pass,
covering what is checked, what problem it solves, and the
implementation details.

%----------------------------------------------------------
\section{Features}

%----------------------------------------------------------
\subsection{Branch Checks}

\textbf{What it does.} A \texttt{refute false} is prepended to
the body of every \texttt{then}-branch and, when an
\texttt{else}-branch exists, to the \texttt{else}-body as well.
Two AST patterns are matched. First, a full
\texttt{if}/\texttt{else}: \texttt{Branch([(cond, thenBody),
(true, elseBody)])} — the second arm's condition is the literal
\texttt{true}, which is the COL encoding of \texttt{else}. Second,
an \texttt{if} without \texttt{else}: \texttt{Branch([(cond,
thenBody)])}. In both cases a check is inserted at the very start
of each arm's body, before any user code runs.

\textbf{What it solves.} When a method's precondition already
implies that a branch condition is always false, the branch body
is dead. The verifier proves everything inside it vacuously but
reports nothing. A common source of this is an overly strong
precondition that rules out certain inputs, or a condition that
contradicts an earlier \texttt{assume}. For example:

\begin{lstlisting}
// After rewriting
requires x > 0;
void f(int x) {
    if (x < 0) {
        refute false; // fires: x>0 /\ x<0 unsat
        x = x + 1;
    } else {
        refute false; // passes: x>0 /\ x>=0 ok
        x = x - 1;
    }
}
\end{lstlisting}

Without the check, the \texttt{then}-branch is silently verified
as dead code. With it, a \texttt{deadBranch} warning points
directly at the branch with the label
\texttt{"then-branch (condition: `x < 0`)"}, telling the developer
exactly which arm is dead and under what condition.

\textbf{Implementation.} The \texttt{dispatch} method on
\texttt{Statement} matches on \texttt{Branch} nodes. For each arm,
\texttt{instrumentBody} is called: it creates a fresh
\texttt{DeadCodeBlame}, builds a
\texttt{Block(Seq(Refute(ff)(blame), dispatch(body)))}, and pushes
the blame onto the \texttt{currentBlame} scoped stack for the
duration of the body rewrite. This means any nested checks
(branches, loops, assumes) created inside the body will see this
blame as their ancestor and can suppress themselves if it fires
(see Section~\ref{sec:cascade}). The condition text is extracted
via \texttt{condText}, which reads the source origin's inline
context; if no source text is available it falls back to
\texttt{<expression>}.

%----------------------------------------------------------
\subsection{Loop Body and Post-Loop Checks}

\textbf{What it does.} Two \texttt{refute false} statements are
inserted for every \texttt{Loop} node: one at the very start of the
loop body, and one immediately after the loop in the surrounding
block. The body check fires when the state $\mathit{inv} \wedge
\mathit{cond}$ is unsatisfiable — the loop is never entered. The
post-loop check fires when $\mathit{inv} \wedge \neg\mathit{cond}$
is unsatisfiable — the loop never terminates, so code after it is
unreachable.

\textbf{What it solves.} A loop body becomes dead when its
condition can never be true given the current state. Less
obviously, code \emph{after} a loop can become dead even without
any branch following it: if the loop invariant combined with the
negated loop condition produces a contradiction, the loop never
exits. This is easy to miss because there is no \texttt{if} or
\texttt{else} that visually signals a dead path. The post-loop
check is the only mechanism that catches this case.

\begin{lstlisting}
loop_invariant inv;
while (cond) {
    refute false; // body: inv /\ cond
    ...
}
refute false;     // post-loop: inv /\ !cond
x = x + 1;       // caught if post-loop fires
\end{lstlisting}

When the loop has a \texttt{LoopInvariant}, both labels include
the invariant text, e.g.\ \texttt{"loop body (condition: `i < n`,
invariant: `i >= 0`)"}, giving the developer full context about
which loop is involved.

\textbf{Implementation.} The \texttt{dispatch} method matches
\texttt{Loop(init, cond, update, contract, body)}. It calls
\texttt{instrumentBody} to produce the instrumented body, then calls
\texttt{makeCheck} a second time to produce the post-loop
\texttt{Refute(false)}. Because \texttt{dispatch} must return a
single \texttt{Statement[Post]}, both the instrumented loop and the
post-loop check are wrapped in a \texttt{Block(Seq(loop, check))}.
\texttt{Block} is itself a \texttt{Statement}, so the return type
is satisfied while both are processed sequentially by the verifier.
The two suppression flags \texttt{invNotEstablishedFired} and
\texttt{invNotMaintainedFired} are described in
Section~\ref{sec:inv-flags}.

%----------------------------------------------------------
\subsection{Switch Case Checks}

\textbf{What it does.} A \texttt{refute false} is inserted after
each \texttt{case} or \texttt{default} label, before any code in
the case body runs. Both \texttt{Case(pattern)} and
\texttt{DefaultCase()} are matched. The label produced for a
\texttt{case} carries the pattern text, e.g.\
\texttt{"switch case (pattern: `3`)"}, while a \texttt{default}
produces \texttt{"switch default case"}.

\textbf{What it solves.} Switch cases can silently become dead when
the switched value is constrained by the method precondition or by
earlier \texttt{assume} statements. For example, if the precondition
states \texttt{x > 0} and the switch contains \texttt{case 0}, that
case is unreachable. The verifier proves its body vacuously without
any warning. The same applies to a \texttt{default} case if earlier
assumptions already cover all possible values.

\textbf{Implementation.} \texttt{dispatch} matches on
\texttt{SwitchCase} nodes. Unlike branches, the check is placed
\emph{after} the case label rather than inside a body, because the
label itself must be reached for the \texttt{switch} to enter the
case. \texttt{makeCheck} is called to produce a
\texttt{DeadCodeBlame} and a \texttt{Refute(false)}, and the result
is emitted as \texttt{Block(Seq(c.rewriteDefault(), blame\_scoped\_check))}.
The \texttt{caseBlame} is pushed onto \texttt{currentBlame} only
around the check statement itself. Because \texttt{SwitchCase}
nodes in COL are plain statement labels — the case body statements
are siblings in the enclosing \texttt{Block}, not children of the
\texttt{SwitchCase} node — there is no nested dispatch scope to
inherit the blame. Cascade suppression therefore does not propagate
from the case-entry check to subsequent case body statements, unlike
branch and loop bodies where \texttt{instrumentBody} wraps the full
body dispatch inside \texttt{currentBlame.having(blame)}.

%----------------------------------------------------------
\subsection{State-Narrowing Statement Checks}

\textbf{What it does.} A \texttt{refute false} is appended after
each statement that unconditionally adds facts to the symbolic
state: \texttt{assume expr}, \texttt{inhale expr}, and
\texttt{lock obj}. However, the two literal-false variants are
treated differently:

\begin{itemize}
  \item \texttt{assume false} / \texttt{inhale false} — the
    statement itself trivially makes the state inconsistent. No
    \texttt{Refute} is inserted; instead a \texttt{DeadBranch}
    warning is emitted immediately and the statement is replaced
    with \texttt{Inhale(false)} (for \texttt{assume}) or left as-is
    (for \texttt{inhale}) so the verifier still sees an
    inconsistent state for downstream checks.
  \item \texttt{assume expr} / \texttt{inhale expr} / \texttt{lock
    obj} — the statement is kept and a \texttt{Refute(false)} is
    appended after it to check whether the strengthened state
    remains consistent.
\end{itemize}

\textbf{What it solves.} These statements narrow the set of
reachable states. If one introduces a contradiction — for example
\texttt{assume x < 0} after a precondition \texttt{x > 0} — all
code following it is dead. Without per-statement checks, only a
later branch or loop body would fire a warning, giving no
indication of which statement was the actual root cause. With the
checks in place, the \texttt{refute false} immediately after the
contradicting statement fires first, directly identifying the
culprit. Subsequent warnings (e.g.\ a dead \texttt{then}-branch
later in the same method) are suppressed by the cascade mechanism.

\begin{lstlisting}
requires x > 0;
void f(int x) {
    assume x < 10;
    refute false;  // (1) passes: 0 < x < 10 ok
    inhale x != 5;
    refute false;  // (2) passes: still consistent
    assume x < 0;  // contradicts x > 0
    refute false;  // (3) fires: root cause here
    if (x == 5) {
        refute false; // (4) suppressed by (3)
    }
}
\end{lstlisting}

Without checks (1)--(3), only (4) would fire and the developer
would not know which \texttt{assume} introduced the contradiction.

\textbf{Implementation.} The helper \texttt{appendCheck} is used
here rather than \texttt{instrumentBody}. The difference is
structural: \texttt{instrumentBody} prepends a check inside a
body, whereas \texttt{appendCheck} emits
\texttt{Block(Seq(stmt.rewriteDefault(), check))}. Crucially, the
blame is pushed onto \texttt{currentBlame} only around the check
statement itself, not around subsequent siblings. This means a
fired check after an \texttt{assume} does not suppress the next
\texttt{if}-branch: both warnings fire, one labelling the root
cause and one labelling the consequence. Note that
\texttt{appendCheck} calls \texttt{node.rewriteDefault()} rather
than \texttt{dispatch}, which is correct for these simple
statements that have no interesting children needing further
instrumentation.

\textbf{Lock statements.} \texttt{lock obj} is treated as a
state-narrowing statement because acquiring a lock transfers
permissions from the lock invariant into the current thread's
state, narrowing the heap. If the lock invariant is inconsistent
with the current state, all code after the lock is dead. The label
produced is \texttt{"code after locking `obj`"}.

%----------------------------------------------------------
\subsection{Parallel and Atomic Block Checks}

\textbf{What it does.} A \texttt{refute false} is prepended to
the body of both \texttt{ParBlock} and \texttt{ParAtomic} nodes.
\texttt{ParBlock} is handled inside the overridden
\texttt{dispatch(ParRegion)} method. \texttt{ParAtomic} is handled
in \texttt{dispatch(Statement)} since it is a \texttt{Statement}
node. In both cases \texttt{instrumentBody} is called on the
block's content.

\textbf{What it solves.} Parallel blocks carry their own
preconditions in the form of permission requirements and guards.
If these cannot be satisfied given the current thread's state —
for instance, if the required resource is already consumed earlier
— the block body is dead. This would not be caught by branch or
loop checks because \texttt{ParBlock} and \texttt{ParAtomic} are
distinct AST nodes with their own dispatch path.
\texttt{ParAtomic} additionally requires that the relevant atomic
invariant holds; if the invariant is inconsistent, the atomic
body is dead.

\textbf{Implementation.} \texttt{dispatch} for \texttt{ParRegion}
is overridden separately from \texttt{dispatch} for
\texttt{Statement} because these are different dispatch targets in
the COL rewriter. For \texttt{ParBlock}, the \texttt{content}
field is replaced with the result of \texttt{instrumentBody(block,
"parallel block body", content)}. For \texttt{ParAtomic}, the
same is done with \texttt{"atomic block body"}. Both calls push a
fresh \texttt{DeadCodeBlame} onto \texttt{currentBlame}, so any
nested checks inside the block body will correctly propagate
suppression. Instrumentation is skipped entirely if
\texttt{methodBlame} has no top (i.e.\ we are outside a method),
matching the same guard applied to all other checks.

%----------------------------------------------------------
\subsection{Intentional Dead Code (\texttt{assert false} Cutoff)}

\textbf{What it does.} When the rewriter processes a
\texttt{Block} and encounters a literal \texttt{assert false}
(matched as \texttt{Assert(BooleanValue(false))}), all statements
\emph{after} it in the same block are dispatched inside an
\texttt{afterAssertFalse} scope rather than normally. The guard at
the top of \texttt{dispatch} checks whether
\texttt{afterAssertFalse} is non-empty; if so, it skips all
instrumentation and calls \texttt{stat.rewriteDefault()} instead.
This is a hard cutoff: no \texttt{refute false} is inserted
anywhere after the \texttt{assert false}, regardless of nesting.

\textbf{What it solves.} Developers sometimes write dead branches
intentionally — for example, to assert that a particular case is
logically impossible and communicate this intent to reviewers. In
Frama-C/WP, the same convention is used: placing
\texttt{assert \textbackslash false} at the start of a branch body
signals that the branch is intentionally dead and suppresses the
smoke-test warning. Our implementation mirrors this exactly.
Without such a mechanism, every dead branch — including deliberate
ones — would generate a warning, creating noise that obscures real
problems.

\begin{lstlisting}
requires x > 0;
void f(int x) {
    if (x < 0) {
        assert false;  // intentional marker
        x = x + 1;    // no refute inserted
        if (x < 2) {  // no refute here either
            x = x + 1;
        }
    }
    if (x > 0) {
        refute false;  // live: checked normally
        x = x - 1;
    }
}
\end{lstlisting}

Note that the second \texttt{if} is outside the
\texttt{assert false} block, so it is still checked normally.
The cutoff applies only to statements in the same \texttt{Block}
after the \texttt{assert false}, and recursively to any nested
blocks via the \texttt{afterAssertFalse} scope.

\textbf{Implementation.} Inside the \texttt{Block} case of
\texttt{dispatch}, \texttt{indexWhere} scans the statement list for
the first \texttt{Assert(BooleanValue(false))}. If none is found
(\texttt{cutoffIdx < 0}), all statements are dispatched normally.
If found, the list is split at \texttt{cutoffIdx + 1}: statements
before and including the \texttt{assert false} are dispatched
normally, and the remaining statements are dispatched inside
\texttt{afterAssertFalse.having(()) \{...\}}. The
\texttt{afterAssertFalse} stack is a \texttt{ScopedStack[Unit]};
its non-emptiness is the sole signal used by the guard.

%----------------------------------------------------------
\subsection{Cascade Suppression}
\label{sec:cascade}

\textbf{What it does.} When a dead-code check fires, all nested
checks inside the same branch body are automatically suppressed.
The suppression works through a chain of blame objects: each
\texttt{DeadCodeBlame} holds a reference to its parent's
\texttt{firedOrSuppressed} flag and checks it before emitting a
warning. If any ancestor has fired, the child sets its own
\texttt{firedOrSuppressed} flag (so its children are also
suppressed) but does not emit a \texttt{deadBranch} warning.

\textbf{What it solves.} If a branch is dead, every statement
inside it is trivially dead too: nested branches, loop bodies, and
post-\texttt{assume} checks all fire vacuously because the
enclosing path condition is already \texttt{false}. Without
suppression, a single dead outer branch would generate a cascade
of warnings for every inner check, making it nearly impossible to
identify the root cause.

\begin{lstlisting}
requires x < 0;
void f(int x) {
    if (x > 0) {
        refute false;  // (1) fires: dead
        if (x > 10) {
            refute false; // (2) suppressed by (1)
        } else {
            refute false; // (3) suppressed by (1)
        }
        assume x == 5;
        refute false;  // (4) suppressed by (1)
    }
}
\end{lstlisting}

Only warning (1) reaches the developer, pointing at the outermost
dead branch. All inner checks are silenced.

\textbf{Implementation.} The \texttt{DeadCodeBlame} class holds
two fields: \texttt{firedOrSuppressed: Boolean} (mutable, private)
and \texttt{ancestorFired: () => Boolean} (a lambda). When
\texttt{blame} is called, it first sets \texttt{firedOrSuppressed
= true} unconditionally — this ensures its own children are
suppressed even if it does not emit a warning — then checks
\texttt{ancestorFired()}. Only if no ancestor has fired does it
forward \texttt{DeadBranch} to \texttt{methodBlame}.

The \texttt{ancestorFired} lambda is constructed in
\texttt{makeCheck}. It reads \texttt{currentBlame.topOption}: if
a parent blame is on the stack, the lambda closes over
\texttt{parentBlame.didFire}; otherwise it returns \texttt{false}.
An optional \texttt{extraSuppressor} lambda is OR-ed in for cases
where suppression depends on external state (e.g.\ loop invariant
flags) rather than the blame stack. This gives a single unified
predicate: \texttt{() => base() || extraSuppressor()}.

Sequential checks at the same nesting level — e.g.\ an
\texttt{assume} check followed by a branch check in the same block
— are not in an ancestor/descendant relationship, so they fire
independently. This is intentional: the developer sees one warning
for the statement that introduced the contradiction and one for
the branch that became dead as a consequence.

%----------------------------------------------------------
\subsection{Loop Invariant Suppression Flags}
\label{sec:inv-flags}

\textbf{What it does.} Two boolean \texttt{var}s are created per
loop: \texttt{invNotEstablishedFired} and
\texttt{invNotMaintainedFired}. They are set by a wrapper placed
around the loop invariant's blame object and are passed as
\texttt{extraSuppressor} lambdas to the body check and post-loop
check respectively.

\textbf{What it solves.} A loop invariant failure can make the
body or post-loop check spuriously fire even though the code is
not actually dead — the invariant simply could not be proved. The
two error kinds cascade differently and require separate treatment:

\begin{itemize}
  \item \texttt{LoopInvariantNotEstablished} — the invariant does
    not hold before the loop is entered. Silicon cannot assume it
    at the top of the loop body, so the body check and the
    post-loop check both fire spuriously. Both must be suppressed.
  \item \texttt{LoopInvariantNotMaintained} — the invariant holds
    at the start of each iteration (Silicon \texttt{assume}s it),
    but the body can falsify it. The body check is therefore not
    spurious — dead-branch findings inside the body remain genuine.
    However, after the failed maintenance assertion Silicon emits
    an implicit \texttt{assume false}, which leaks into the
    post-loop state. Only the post-loop check is suppressed.
\end{itemize}

\begin{lstlisting}[language=Scala]
val bodySuppress: () => Boolean =
  () => invNotEstablishedFired
val postLoopSuppress: () => Boolean =
  () => invNotEstablishedFired || invNotMaintainedFired
\end{lstlisting}

\textbf{Implementation.} When the loop's contract is a
\texttt{LoopInvariant}, the pass replaces its blame with a
\texttt{Blame[LoopInvariantFailure]} wrapper. The wrapper inspects
the error variant on each call: a
\texttt{LoopInvariantNotEstablished} sets
\texttt{invNotEstablishedFired = true}; a
\texttt{LoopInvariantNotMaintained} sets
\texttt{invNotMaintainedFired = true}. The original blame is always
forwarded so the user still sees the invariant error. The flags are
then closed over by the \texttt{bodySuppress} and
\texttt{postLoopSuppress} lambdas, which are passed as
\texttt{extraSuppressor} arguments to \texttt{makeCheck}. Because
the flags are local \texttt{var}s defined inside the
\texttt{Loop} case, they are per-loop and never interfere across
different loops in the same method.

%----------------------------------------------------------
\subsection{Loop Invariant Satisfiability Check}

\textbf{What it does.} A separate pass,
\texttt{CheckLoopInvariantSatisfiability}, runs after
\texttt{DetectDeadCode}. For every \texttt{Loop} whose contract is
a \texttt{LoopInvariant(inv)}, it generates an isolated proof
context that inhales \texttt{inv} and immediately runs
\texttt{refute false}. If the invariant is internally contradictory
(e.g.\ \texttt{x > 0 \&\& x < 0}), inhaling it makes
\texttt{false} derivable, the refute fires, and a
\texttt{LoopInvariantUnsatisfiable} error is reported on the
invariant node.

\textbf{What it solves.} Normal loop verification cannot detect an
unsatisfiable invariant in two important situations. First, when
the method precondition is \texttt{false}, loop verification is
vacuously skipped — the loop is never reached, so the invariant is
never checked:

\begin{lstlisting}
requires false;
void f(int x) {
    // invariant is broken but never checked
    loop_invariant x > 0 && x < 0;
    while (x > 0) { x--; }
}
\end{lstlisting}

Second, even without a false precondition, a contradictory
invariant currently produces \texttt{invariantNotEstablished} —
because the initial state does not satisfy the invariant. This
error is ambiguous: the invariant might simply not hold at loop
entry, or it could be fundamentally broken. After the
satisfiability check, \texttt{invariantUnsatisfiable} fires first,
giving a clear, unambiguous diagnosis. The pass also catches
resource invariants where no state can hold the required
permissions, e.g.\ \texttt{Perm(c, write) ** Perm(c, write)}:
inhaling this in an isolated context immediately produces an
inconsistent heap, triggering the refute.

\begin{lstlisting}[language=Scala]
// For: loop_invariant inv; while (cond) { ... }
Block(Seq(
  Extract(
    FramedProof(
      pre  = tt,
      body = Block(Seq(
        Inhale(dispatch(inv)), // assume invariant holds
        Refute(ff)(blame),     // unsat? -> fires
      )),
      post = tt,
    )(TrueSatisfiable),
    decreases = None,
  )(...),
  loop.rewriteDefault(),  // original loop follows
))
\end{lstlisting}

\textbf{Implementation.} \texttt{FramedProof(true, body, true)}
creates an isolated proof context with no preconditions and no
postconditions. \texttt{Extract} wraps it so that
\texttt{EncodeExtract} (which runs later in the pipeline) lifts it
into a standalone generated procedure with signature
\texttt{requires true; ensures true;}. This gives complete heap
isolation: the check is independent of the caller's state, the
method's precondition, and any other loop in the same method.
Inhaling \texttt{dispatch(inv)} inside this context makes an
unsatisfiable invariant produce an immediately inconsistent heap
state, causing \texttt{Refute(false)} to fire. The blame converts
the \texttt{RefuteFailed} into
\texttt{LoopInvariantUnsatisfiable(li)}, pointing the error at the
\texttt{LoopInvariant} node in the AST. The pass must run before
\texttt{EncodeExtract} (which encodes the \texttt{Extract} node)
and before \texttt{EncodeProofHelpers} (which encodes the
\texttt{FramedProof} node). The generated \texttt{Refute} nodes
are not transformed by \texttt{RefuteToInvertedAssert}; they are
translated directly to Viper's native \texttt{refute} statement in
\texttt{ColToSilver} via the \texttt{RefutePlugin}.

%----------------------------------------------------------
\section{Possible Improvements}

\textbf{Requires-false detection.} A \texttt{refute false} at
method entry would detect a false precondition directly. This is
not currently implemented because \texttt{CheckContractSatisfiability}
already covers this case; adding it inside \texttt{DetectDeadCode}
would produce a redundant check and potentially a duplicate warning.

\textbf{Sequential cascade suppression.} When \texttt{assume false}
fires, a subsequent \texttt{if} branch fires too as an independent
warning, since neither is an ancestor of the other in the blame
tree. Both warnings are correct — one for the root cause, one for
the consequence — but when many statements follow a single
contradictory \texttt{assume}, the output can be noisy. A fix
would be to track per-\texttt{Block} shared state: each statement's
\texttt{appendCheck} result would be visible to subsequent siblings,
allowing them to suppress themselves when a predecessor already
fired.

\end{document}
